% ══════════════════════════════════════════════════════════════
\section{CENA 2 --- O Chip que Ferve: Integral Dupla \hfill \normalfont\small\color{corTempo}2:00 -- 10:00}
% ══════════════════════════════════════════════════════════════

\subsection{2.1 -- O problema físico}
\tempo{2:00 -- 3:30}

\begin{visual}
\faFilm\ Câmera voa sobre a superfície do chip como drone. Cores indicam temperatura:
azul nas bordas, laranja/vermelho no centro.
\end{visual}

\begin{fala}
Imagine que você precisa saber a temperatura real de um processador.
Parece fácil, né? Coloca um sensor e pronto.
\end{fala}

\begin{visual}
\faFilm\ Um único ponto de medição aparece no chip, mostrando 72 {\textdegree}C.
Mas outras regiões mostram 45 {\textdegree}C, 160 {\textdegree}C, 55 {\textdegree}C simultaneamente.
O sensor pega apenas um deles.
\end{visual}

\begin{fala}
O problema é que um sensor pontual te dá \textbf{um} valor. Mas um chip real
tem pontos a 45 graus, regiões a 100 graus, e hotspots a 160 graus --
tudo ao mesmo tempo. Um único número não conta a história inteira.
\end{fala}

\begin{fala}
Para enxergar o chip de verdade, a gente precisa de uma ferramenta que
\textbf{integre} toda a superfície. E é exatamente aí que entra a integral dupla.
\end{fala}

\begin{tela}
\faDesktop\ \textbf{Equação aparece na tela, termo por termo:}
\end{tela}

\begin{mathblock}
\textbf{\color{corMath}\faSuperscript\ A fórmula da temperatura média:}

\vspace{6pt}
\[
  T_{\text{média}}
  \;=\;
  \frac{1}{\underbrace{\text{Área}}_{\text{tamanho do chip}}}
  \iint_{A}
  \underbrace{T(x,\,y)}_{\text{temp.\ em cada ponto}}
  \,\mathrm{d}x\,\mathrm{d}y
\]

\vspace{4pt}
{\small\color{gray}
\textbf{Como ler isso:} varremos cada ponto $(x,y)$ da superfície, lemos a temperatura
lá, somamos tudo e dividimos pela área total. É a média de verdade -- não um chute.}
\end{mathblock}

\begin{fala}
Geometricamente, a integral dupla está calculando o \textbf{volume} debaixo
da superfície de temperatura. Dividindo pela área da base, obtemos a
\textbf{altura média} -- que é exatamente a temperatura média que queremos.
\end{fala}

\begin{visual}
\faFilm\ Animação 3D: superfície T(x,y) flutua sobre o chip. Uma fatia horizontal
mostra a altura média. O volume entre superfície e chip fica colorido de laranja.
\end{visual}

% ─────────────────────────────────────────────────────────────
\subsection{2.2 -- Escolhendo a função T(x,y)}
\tempo{3:30 -- 5:00}
% ─────────────────────────────────────────────────────────────

\begin{fala}
Mas que função a gente usa pra representar o calor do chip?
Hotspots no mundo real têm um formato muito característico:
picos concentrados -- quentes no centro, esfriando em volta.
Sabe quem tem esse formato? A Gaussiana.
\end{fala}

\begin{visual}
\faFilm\ Animação: curva de sino em 2D se expande para montanha 3D centralizada no chip.
Legenda: ``Gaussiana = pico natural do mundo físico''.
\end{visual}

\begin{tela}
\faDesktop\ \textbf{Função na tela:}
\end{tela}

\begin{mathblock}
\textbf{\color{corMath}\faSuperscript\ Função temperatura escolhida:}

\[
  T(x,\,y)
  \;=\;
  \underbrace{60}_{\substack{\text{temperatura}\\\text{de base ({\textdegree}C)}}}
  \;+\;
  \underbrace{100}_{\substack{\text{pico do}\\\text{hotspot ({\textdegree}C)}}}
  \cdot\;
  \exp\!\Bigl[
    -\,\frac{(x-0{,}5)^2 + (y-0{,}5)^2}{2\times(0{,}1)^2}
  \Bigr]
\]

\vspace{6pt}
{\small\color{gray}
Domínio: $[0,1]\times[0,1]$ (chip de $1\,\text{cm} \times 1\,\text{cm}$) \quad
Centro do hotspot: $(0{,}5;\;0{,}5)$ \quad
$\sigma = 0{,}1\,\text{cm}$ \quad
$T_{\max} = 160\,\text{\textdegree C}$}
\end{mathblock}

\ferramenta{Matplotlib -- heatmap da função T(x,y)}

\begin{lstlisting}[caption={Mapa de calor com Matplotlib}]
import numpy as np
import matplotlib.pyplot as plt

# Discretização (os mesmos dados do roteiro)
x = np.linspace(0, 1, 100)
y = np.linspace(0, 1, 100)
X, Y = np.meshgrid(x, y)
T = 60 + 100 * np.exp(-((X-0.5)**2 + (Y-0.5)**2) / 0.02)

fig = plt.figure(figsize=(10, 7))
ax = fig.add_subplot(111, projection='3d')
surf = ax.plot_surface(X, Y, T, cmap='hot', edgecolor='none', alpha=0.9)
ax.set_title('Hotspot Gaussiano no Processador')
ax.set_zlim(60, 160)
plt.colorbar(surf, label='Temperatura (°C)')
plt.show()
\end{lstlisting}

\begin{visual}
\faFilm\ Heatmap aparece na tela. Vermelho intenso no centro, azul nas bordas.
Seta apontando para o centro: ``HOTSPOT: 160 {\textdegree}C''.
\end{visual}

% ─────────────────────────────────────────────────────────────
\subsection{2.3 -- Calculando a integral dupla}
\tempo{5:00 -- 8:00}
% ─────────────────────────────────────────────────────────────

\begin{fala}
Agora sim. Vamos calcular a temperatura média de verdade.
Substituindo nossa função na fórmula da integral dupla:
\end{fala}

\begin{mathblock}
\textbf{\color{corMath}\faSuperscript\ Passo 1 -- Escrever a integral:}

\[
  T_{\text{média}}
  = \int_0^1\!\!\int_0^1
    \Bigl[\,60 + 100\,e^{-\frac{(x-0{,}5)^2+(y-0{,}5)^2}{0{,}02}}\,\Bigr]
    \,\mathrm{d}x\,\mathrm{d}y
\]

\textbf{\color{corMath}Passo 2 -- Separar em dois pedaços:}

\[
  = \underbrace{\int_0^1\!\!\int_0^1 60\,\mathrm{d}x\,\mathrm{d}y}_{=\;60}
  \;+\;
  100
  \cdot
  \underbrace{\int_0^1 e^{-\frac{(x-0{,}5)^2}{0{,}02}}\,\mathrm{d}x}_{\mathcal{I}}
  \cdot
  \underbrace{\int_0^1 e^{-\frac{(y-0{,}5)^2}{0{,}02}}\,\mathrm{d}y}_{=\;\mathcal{I}}
\]

{\small\color{gray}\textbf{Por quê separar?} O exponencial é produto de uma função só
de $x$ por uma só de $y$. Integral separável = produto das integrais. Cada uma vira
a mesma $\mathcal{I}$ -- menos conta!}
\end{mathblock}

\begin{fala}
A parte difícil é calcular $\mathcal{I}$. Ela tem um truque clássico:
mudança de variável.
\end{fala}

\begin{mathblock}
\textbf{\color{corMath}\faSuperscript\ Passo 3 -- Calcular $\mathcal{I}$:}

\begin{align*}
  \mathcal{I}
  &= \int_0^1 e^{-(x-0{,}5)^2/0{,}02}\,\mathrm{d}x \\[6pt]
  &\overset{u\,=\,x-0{,}5}{=}
    \int_{-0{,}5}^{0{,}5} e^{-u^2/0{,}02}\,\mathrm{d}u \\[6pt]
  &\overset{t\,=\,u/(0{,}1\sqrt{2})}{=}
    0{,}1\sqrt{2}
    \int_{-3{,}54}^{3{,}54} e^{-t^2}\,\mathrm{d}t
    \;\approx\;
    0{,}1\sqrt{2}\cdot\sqrt{\pi}
    = 0{,}1\sqrt{2\pi}
\end{align*}

{\small\color{gray}$\text{erf}(3{,}54) \approx 1{,}0000$ -- podemos estender para
$\pm\infty$ sem remorso matemático.}

\[
  \mathcal{I} \approx 0{,}1 \times 2{,}507 = \mathbf{0{,}2507}
\]

\textbf{\color{corMath}Passo 4 -- Resultado final:}

\[
  T_{\text{média}}
  = 60 + 100 \cdot (0{,}2507)^2
  = 60 + 100 \times 0{,}06285
  \qquad
  \boxed{\;T_{\text{média}} \approx 66{,}28\,\text{\textdegree C}\;}
\]
\end{mathblock}

\begin{visual}
\faFilm\ O número 66,28 {\textdegree}C aparece em destaque no heatmap.
Seta para o pico: ``HOTSPOT: 160 {\textdegree}C''. Seta para a média: ``MÉDIA: 66,28 {\textdegree}C''.
Diferença de 94 {\textdegree}C destacada em vermelho.
\end{visual}

\ferramenta{WolframAlpha -- validação da integral dupla}
\begin{lstlisting}[language=bash]
integrate (60 + 100 * exp(-((x-0.5)^2 + (y-0.5)^2)/0.02)) dx dy from x=0 to 1, y=0 to 1
\end{lstlisting}

\ferramenta{SciPy -- verificação numérica}
\begin{lstlisting}[caption={Integral dupla com SciPy}, language=Python]
from scipy.integrate import dblquad
import numpy as np

def T(y, x):
    return 60 + 100*np.exp(-((x-0.5)**2 + (y-0.5)**2) / 0.02)

resultado, erro = dblquad(T, 0, 1, lambda x: 0, lambda x: 1)
print(f"T_media = {resultado:.4f} C")
print(f"Erro    = {erro:.2e}")
\end{lstlisting}

% ─────────────────────────────────────────────────────────────
\subsection{2.4 -- O que esse número significa?}
\tempo{8:00 -- 10:00}
% ─────────────────────────────────────────────────────────────

\begin{fala}
Sessenta e seis vírgula vinte e oito graus. Parece OK, né?
Mas espera -- o hotspot estava em 160 graus!
\end{fala}

\begin{fala}
A temperatura média mascara o perigo. A \textbf{temperatura de junção}
de um processador típico fica entre 100 e 125 graus.
Nosso hotspot ultrapassa isso facilmente.
\end{fala}

\begin{visual}
\faFilm\ Termômetro animado com três marcas coloridas:
\textcolor{ForestGreen}{66,28 {\textdegree}C -- média (OK)} /
\textcolor{Orange}{100 {\textdegree}C -- limite} /
\textcolor{corAlerta}{160 {\textdegree}C -- hotspot (PERIGO!)}.
Sirene animada pisca ao lado do 160 {\textdegree}C.
\end{visual}

\begin{fala}
Sem a integral dupla, um sensor único poderia estar medindo qualquer
um desses valores e nos dando uma visão completamente errada.
A integral dupla é o \textbf{raio-X da pastilha de silício}.
\end{fala}

\begin{tela}
\faDesktop\
\begin{center}
\large\itshape
``A integral dupla não é exercício de cálculo.\\
É diagnóstico. É prevenção.''
\end{center}
\end{tela}